\section*{الفصل \ref{ch:b2:deep-time-climate} --- المناخ عبر الزمن الجيولوجي}

\begin{solution}{exo:b2:deep-time-climate:1}
(أ) $900 + 37100 + 700 + 3 = \qty{38703}{PgC}$، أي نحو 65 مرّة قدر \qty{591}{PgC} في الهواء.
(ب) $\num{3.5e6}\times0.012011 = \qty{42000}{PgC}$. (ج) $591/42000 = \qty{1.4}{\percent}$.
\end{solution}

\begin{solution}{exo:b2:deep-time-climate:2}
(أ) المدخلات $11.3 + 3 + 12.3 + 8.1 + 2.0 = \qty{36.7}{Tmol.yr^{-1}}$؛ والمخرجات $26 + 16 =
\qty{42}{Tmol.yr^{-1}}$، وهما متساويان ضمن حدود الشكّ فيهما. (ب) $(11.3 + 3)/36.7 = 0.39$: أي دون
النصف بكثير؛ فمعظم المدخلات كربون رسوبي قديم. (ج) $\num{3.5e6}/42 = \qty{8.3e4}{yr}$.
\end{solution}

\begin{solution}{exo:b2:deep-time-climate:3}
(أ) \ce{Mg2SiO4 + 4 CO2 + 2 H2O -> 2 Mg^{2+} + 4 HCO3^- + SiO2}. (ب) مع \ce{2 Mg^{2+} + 4 HCO3^-
-> 2 MgCO3 + 2 CO2 + 2 H2O}، يكون التفاعل الصافي \ce{Mg2SiO4 + 2 CO2 -> 2 MgCO3 + SiO2}: أي مولان
من ثنائي أكسيد الكربون لكلّ مول من الفورستريت. (ج) $\qty{1}{PgC} = 10^{15}/12.011 =
\qty{8.33e13}{mol}$ من الكربون، تحتاج إلى \qty{4.16e13}{mol} من الفورستريت، أي $4.16\times10^{13}
\times\qty{140.7}{g} = \qty{5.86}{Pg}$، قرابة ستة مليارات طنّ.
\end{solution}

\begin{solution}{exo:b2:deep-time-climate:4}
التجلّد الهيرنانتي، من 445 إلى \qty{443}{Ma}، حالة جليدية؛ وأقصى جليد في أواخر الباليوزوي، من 307 إلى
\qty{299}{Ma}، حالة جليدية؛ والتوروني، من 94 إلى \qty{89}{Ma}، حالة دفيئة (الأدفأ في دهر الحياة الظاهرة)؛
والحدّ الأقصى الحراري PETM، من 55.9 إلى \qty{55.7}{Ma}، حالة دفيئة؛ والأمثل المناخي EECO، من 53 إلى \qty{49}{Ma}، حالة دفيئة؛ وبداية
الغطاء الجليدي القطبي الجنوبي، نحو \qty{34}{Ma}، بداية الحالة الجليدية الحالية؛ والأغطية الجليدية الشمالية،
\qty{2.7}{Ma}، حالة جليدية.
\end{solution}

\begin{solution}{exo:b2:deep-time-climate:5}
(أ) $63000/(8.314\times288^2) = \qty{0.0914}{K^{-1}}$، وهي تُقرَّب إلى 0.09. (ب) $74000/(8.314
\times288^2) = \qty{0.107}{K^{-1}}$. (ج) $\ln1.2/0.09 = \qty{2.03}{K}$ و$\ln 1.2/0.107 =
\qty{1.70}{K}$. (د) في الحالة المستقرّة يجب أن تساوي التجوية انطلاق الغازات، وهذا يثبّت $e^{k
\Delta T}$ ومن ثمّ $\Delta T$؛ أمّا الحساسية فلا تحدّد إلا كمّية ثنائي أكسيد الكربون التي تُحدث ذلك
الاحترار.
\end{solution}

\begin{solution}{exo:b2:deep-time-climate:6}
هنا $k_T = 0.0914$. (أ) الدقيق $\exp[7577.5\times(1/288 - 1/293)] = 1.567$؛ والتقريب
$e^{0.0914\times5} = 1.579$: أعلى بنسبة \qty{0.8}{\percent}. (ب) الدقيق $e^{2.482} = 11.97$؛
والتقريب $e^{2.741} = 15.50$: أعلى بنسبة \qty{30}{\percent}. (ج) $\ln(\text{النسبة}) = (E_{\text
a}/R)[x/T_0^2 - x/(T_0(T_0 + x))] = (E_{\text a}/R)\,x^2/(T_0^2(T_0 + x))$. وبجعله مساويًا للقيمة $\ln1.1$:
$x^2 = 1.043\,(288 + x)$، فيكون $x = \qty{17.9}{K}$. وبعد نحو \qty{18}{K} من الاحترار تبالغ
الصيغة الأسّية في تقدير التجوية بأكثر من العُشر.
\end{solution}

\begin{solution}{exo:b2:deep-time-climate:7}
(أ) $1444/276 = 5.23$، و$3.93\log_2 5.23 = \qty{9.4}{W.m^{-2}}$. (ب) تساوي تقريبًا تمامًا
القيمة \qty{9}{W.m^{-2}} للشمس: فالتأثيران يتعادلان. (ج) $(1444 - 276)/419.5 =
\qty{2.8}{ppm}$ لكلّ مليون سنة. أمّا القيمة \qty{3.4}{ppm} فتأتي من خطّ مستقيم ملائم لكلّ
المعطيات، يزن كلّ تقدير بدلًا من طرفي المنحنى المنعّم.
\end{solution}

\begin{solution}{exo:b2:deep-time-climate:8}
(أ) من $3000/20000 = \qty{0.15}{PgC.yr^{-1}}$ إلى $7000/3000 = \qty{2.3}{PgC.yr^{-1}}$، وهو مدى
يحتوي المدى 0.3 إلى 1.5 الذي تذكره الهيئة الحكومية الدولية. (ب) $35\times12.011/44.01 = \qty{9.6}{PgC.yr^{-1}}$، أي من 6 إلى
32 مرّة قدر معدّل الحدّ الأقصى الحراري البالغ 0.3 إلى \qty{1.5}{PgC.yr^{-1}}. (ج) $3.93\log_2(2000/900) = \qty{4.5}{W.m^{-2}}$،
و$3\log_2 2.22 = \qty{3.5}{\celsius}$: أي نحو سبعين في المئة من \qty{5}{\celsius}. ولا بدّ أن تعوّض الباقي تغذياتٌ
راجعة أخرى أو مصدر آخر للاحترار، كالميثان.
\end{solution}

\begin{solution}{exo:b2:deep-time-climate:9}
(أ) الآن $1.2\,W_0e^{k\Delta T} = V = W_0$، فيكون $\Delta T = -\ln 1.2/0.09 = \qty{-2.03}{K}$.
(ب) $278\times2^{-2.03/3} = \qty{174}{ppm}$. (ج) لا: ففي أيّ حالة مستقرّة تساوي التجوية
انطلاق الغازات، الذي لم يتغيّر. وقابلية التجوية الأعلى يعوّضها مناخ أبرد.
\end{solution}

\begin{solution}{exo:b2:deep-time-climate:10}
(أ) $k = 0.119$: $\tau = 591\times0.693/(0.08\times0.119\times3\times0.138) = \qty{104}{kyr}$.
(ب) $S = 8$: $\tau = \qty{51.5}{kyr}$. (ج) كلاهما يجعل منظّم الحرارة أسرع، لأن $\tau \propto
1/(kS)$. والشكّ في $E_{\text a}$ بعامل اثنين تقريبًا، وفي $S$ على هذه المقاييس الزمنية بعامل
مماثل (3 مقابل 5 إلى \qty{8}{K}): فكلاهما مهمّ بالقدر نفسه تقريبًا.
\end{solution}

\begin{solution}{exo:b2:deep-time-climate:11}
(أ) يبقى $0.08\times5000 = \qty{400}{PgC}$ في الهواء: $C/C_0 = 991/591 = 1.677$، أي
\qty{466}{ppm}؛ و$\Delta T = 3\log_2 1.677 = \qty{2.24}{K}$. (ب) $W - W_0 = 0.138\,(e^{0.09\times2.24}
- 1) = \qty{0.0308}{PgC.yr^{-1}}$، أي تجوية أكثر بنسبة \qty{22}{\percent}؛ و$5000/0.0308 =
\qty{162}{kyr}$. (ج) $\tau = \qty{137}{kyr}$ و$\tau\ln10 = \qty{315}{kyr}$. (د) الاحترار لا يزداد
إلا كلوغاريتم ثنائي أكسيد الكربون، فيرفع فائضٌ كبير التجويةَ بأقلّ ممّا يفترضه
القانون الخطّي؛ فيبدأ النزع أبطأ.
\end{solution}

\begin{solution}{exo:b2:deep-time-climate:12}
(أ) قرب $M^*$، $\mathrm d\Delta T/\mathrm dM = Sf/(A_0\ln 2\,(1 + fM^*/A_0)) = Sf/(A_0\ln 2\,
C^*/C_0)$ و$\mathrm dW/\mathrm dM = kW\,\mathrm d\Delta T/\mathrm dM$، مع $W = V$ هناك؛
والمقلوب هو $\tau'$. (ب) $C^*/C_0 = 2^{2.03/3} = 1.597$: $\tau' = 137.4\times1.597/1.2 =
\qty{183}{kyr}$. (ج) أبطأ: فعند تركيز مرتفع لثنائي أكسيد الكربون يُحدث كلّ بيتاغرام مضاف من الكربون احترارًا أقلّ،
لأن التأثير لوغاريتمي، فتستجيب له التجوية أقلّ.
\end{solution}

\begin{solution}{pb:b2:deep-time-climate:1}
\textbf{1.} $1.3 + 1.3 + 4.5 + 4.2 = \qty{11.3}{Tmol.yr^{-1}}$، أي $\qty{0.136}{PgC.yr^{-1}}$.
\textbf{2.} $11.5\times0.012011 = \qty{0.138}{PgC.yr^{-1}}$.
\textbf{3.} يتّفقان في حدود نحو \qty{2}{\percent}: فتجوية السيليكات توازن انطلاق الغازات
البركاني، كما تقتضي \omterm{def:b2:deep-time-climate:carbon-cycle}{الدورة الطويلة الأمد للكربون}.
\textbf{4.} $\num{5.1e18}/0.02897 = \qty{1.76e20}{mol}$ من الهواء؛ و$\num{591e15}/12.011 =
\qty{4.92e16}{mol}$ من الكربون؛ و$\num{4.92e16}/\num{1.76e20} = \num{2.80e-4}$، أي \qty{280}{ppm}.
\textbf{5.} $591/0.136 = \qty{4350}{yr}$.
\textbf{6.} $1 + 0.08 = 1.08$: أي أكثر بنسبة \qty{8}{\percent} فقط.
\textbf{7.} $63000/(8.314\times288^2) = \qty{0.0914}{K^{-1}}$، وهي تُقرَّب إلى 0.09.
\textbf{8.} $e^{0.09} = 1.09$: أي أكثر بنسبة \qty{9}{\percent} لكلّ كلفن.
\textbf{9.} $\Delta T = S\log_2(1 + fM/A_0)$؛ ومشتقّه عند 0 هو $Sf/(A_0\ln2) =
3\times0.08/(591\times0.693) = \qty{5.86e-4}{K.PgC^{-1}}$.
\textbf{10.} $W_0k\times\num{5.86e-4} = 0.138\times0.09\times\num{5.86e-4} = \qty{7.28e-6}{yr^{-1}}$.
\textbf{11.} مقلوبه هو $A_0\ln2/(fkSW_0)$، أي $\tau$ في القضية، لأن
المعادلة الخطّية هي $\mathrm dM/\mathrm dt = -(\mathrm dW/\mathrm dM)\,M$.
\textbf{12.} الأربع جميعًا: فالثابت $\tau$ يتناسب عكسيًا مع كلّ منها.
\textbf{13.} $1/\num{7.28e-6} = \qty{1.37e5}{yr}$.
\textbf{14.} $137\times0.09/0.119 = \qty{104}{kyr}$.
\textbf{15.} $k_T = 0.107$: $137\times0.09/0.107 = \qty{115}{kyr}$.
\textbf{16.} أقصر بعامل من 1.5 إلى 7: فالنموذج ذو الخزّان الواحد وكسره الثابت $f$ يجعلان
منظّم الحرارة أسرع ممّا ينبغي.
\textbf{17.} $\num{50e6}/\num{1.37e5} = 364$: فقد كان لمنظّم الحرارة وقت ليوازن الدورة في كلّ
مرحلة، فتبع الانخفاض تغيّرات بطيئة في انطلاق الغازات أو في قابلية التجوية.
\textbf{18.} $16\times0.137 = \qty{2.2}{ppm}$.
\textbf{19.} $0.08\times5000 = \qty{400}{PgC}$؛ و$278\times991/591 = \qty{466}{ppm}$.
\textbf{20.} $3\log_2 1.677 = \qty{2.24}{K}$.
\textbf{21.} $0.138\,(e^{0.202} - 1) = \qty{0.0308}{PgC.yr^{-1}}$، أي \qty{22.3}{\percent}.
\textbf{22.} $137\times\ln10 = \qty{315}{kyr}$.
\textbf{23.} ربّما كانت التجوية أشدّ حساسيةً (الجريان السطحي، وطاقة تنشيط أكبر،
والتغذيات الراجعة البطيئة لحساسية نظام الأرض الأكبر)؛ وربّما عاد الكربون العضوي على اليابسة وفي البحر إلى
النموّ ونزع الكربون من الهواء أسرع من تجوية السيليكات.
\textbf{24.} في هذا النموذج ينزع \omterm{def:b2:deep-time-climate:thermostat}{منظّم الحرارة بتجوية السيليكات} فائضًا من الكربون بزمن
اضمحلال أسّي يقارب \textbf{\qty[text-series-to-math=true]{137}{} ألف سنة}.
\end{solution}
